\chapter{China's Scientific Machine: AI for Science as Industrial Policy}
\chaptermark{China's Scientific Machine}
\label{ch:china-ai4s}

Chapter~\ref{ch:prize} argued that the largest value in artificial intelligence is scientific, and closed by observing that the contest over that value is genuinely joined. This chapter takes the observation apart on one side of it. Its claim is not that China has won AI for science---the failure ledger of Chapter~\ref{ch:prize} applies to Chinese results exactly as it applies to Western ones, and this chapter applies it, sometimes to results that its own sources report as triumphs. The claim is narrower and harder to dismiss: that China has been organizing for this particular prize longer, more explicitly, and more \emph{industrially} than any Western actor, that the organizing is legible in the primary planning documents rather than inferred from behaviour, and that the shape it has taken---discovery chained by policy text to engineering and to product, running on a domestic substrate, distributed through open artifacts---is the playbook of Chapter~\ref{ch:playbook} applied to the research enterprise itself. The chapter also carries its own best refutation, and carries it first, because the refutation comes from reading the Chinese planning documents precisely rather than enthusiastically, and precision here cuts against the thesis.

\section{The policy record, read precisely}
\label{sec:china-ai4s-policy}

The State Council's \emph{Opinions on Deepening the Implementation of the ``AI+'' Initiative}---Guofa [2025] No.~11, signed 21~August 2025---set out six key actions, and ``AI+ science and technology'' is the first of them, ahead of industry, consumption, welfare, governance and global cooperation~\cite{ai-plus-science} [P]. Ordering in a State Council Opinions document is not decorative; it is the sequence in which ministries are expected to read their own obligations, and placing science ahead of industry in a document whose entire register is industrial is a deliberate statement about where the leadership believes the compounding lies. The science pillar has three sub-clauses. The first directs the acceleration of AI-driven new research paradigms and of ``from 0 to 1'' discovery. The second directs the construction of \emph{scientific large models} and of AI-upgraded research platforms---an instruction to build artifacts, not to fund projects. The third is the analytically load-bearing one, and it is why this chapter is titled industrial policy rather than science policy: it directs the integrated development of ``AI-driven technology R\&D, engineering realisation and product landing.'' The document does not treat discovery as an end. It treats discovery as the first stage of a pipeline whose last stage is a product, and it says so in the text rather than leaving the inference to a foreign analyst.

Set against that, the same document's numeric commitments are the chapter's first piece of counter-evidence, and it is a substantial one. The targets in Guofa [2025] No.~11 are \emph{diffusion} targets---penetration of new-generation intelligent terminals and agents above 70\% by 2027 and above 90\% by 2030---and there is \textbf{no numeric target for science at all}~\cite{ai-plus-science} [P]. A planning system whose distinctive competence is hitting the things it measures has, in its flagship AI document, measured adoption and declined to measure discovery. The strong reading of the ordering and the honest reading of the target list point in different directions, and both belong in the record.

The lineage runs back further, and its form matters as much as its content. The Ministry of Science and Technology and the National Natural Science Foundation launched the ``AI for Science'' special deployment on 27~March 2023, naming four basic disciplines---mathematics, physics, chemistry, astronomy---and four applied priorities---drug development, gene research, agricultural breeding, and new materials---together with commitments to domain-specific platforms, to the national next-generation AI public computing platform, to heterogeneous integration of high-performance and intelligent computing centres, and to scientific data sharing~\cite{most-ai4s} [P]. It is a ministerial \emph{deployment of work}, not a numbered regulation with an implementing budget attached, and Western sources that describe it as a ``plan'' overstate its form; this book does not, and the correction is not pedantry, because a deployment of work and a State Council plan carry very different degrees of compulsion on the agencies below them~\cite{china-ai4s}. The traceable competitive instrument that followed is the National Key R\&D Programme's ``Disruptive Technology Innovation'' special programme, whose 2025 AI-for-science call of 13~May 2025 names five directions: frontier science, life and health, new materials, engineering design, and other AI-driven research-system technologies~\cite{most-ai4s}.

Now the precision point that is worth the whole section. The Party's \emph{Recommendations} for the 15th Five-Year Plan, adopted at the Fourth Plenum and released on 28~October 2025, contain in section~14 the operative instruction to \emph{use artificial intelligence to lead a transformation of the scientific research paradigm}, and the official interpretation built around that phrase carries commentary from the CAS academician who founded the field's principal Chinese institute~\cite{fyp15-ai4s} [P]. \textbf{That phrase does not survive into the Outline adopted by the National People's Congress in March 2026.} In the Outline, AI appears under compute, data and diffusion---Chapter~12 on compute, algorithm and data supply, Chapter~13 on digital empowerment---while basic research appears in Chapter~8 \emph{without} AI attached to it~\cite{fyp15-ai4s} [P]. The consequence for anyone citing the 15th Five-Year Plan in this argument is direct: a claim that the plan makes AI for science a central pillar of Chinese \emph{state programming} is a claim about the Party recommendations, not about the plan, and the two documents do different work. The Recommendations state ambition; the Outline allocates. Two readings of the gap are available and this book prints both. The deflationary reading is that the Outline is a resource-allocation instrument and a paradigm shift is not a line item, so the omission means nothing. The consequential reading is that the drafting process, between October and March, converted a statement about \emph{how research should be done} into a set of provisions about \emph{what infrastructure should be bought}---and that the conversion is itself the finding, because it points at exactly the same revealed preference that Section~\ref{sec:china-ai4s-infra} finds in the state's own prize list. Either way, the gap between Party ambition and state programming is real, it is documentable from primary sources, and it is the strongest single piece of evidence against this chapter's thesis.

\begin{casestudy}{Two documents, one plan: how to read the 15th Five-Year Plan on AI for science}
The Five-Year Plan is not a document; it is a sequence of documents with different authors, different force and different functions, and conflating them is the most common error in Western writing about Chinese technology policy. Three artifacts matter here. The \emph{Recommendations} (\emph{jianyi}) are a Party document, adopted at a plenum, stating direction and priority; they are where ``use artificial intelligence to lead a transformation of the scientific research paradigm'' appears, in section~14, in October 2025. The \emph{Outline} (\emph{gangyao}) is a state document, adopted by the National People's Congress the following March, which converts direction into chapters, programmes and indicators; it is the instrument ministries and provinces are actually assessed against. And beneath both sit the sectoral and provincial implementation plans, which is where money appears. The paradigm-transformation phrase lives in the first and does not appear in the second. What \emph{does} survive into the Outline is compute: ultra-large-scale intelligent computing clusters under study and construction, government purchase and leasing of compute, universal accessibility and lower costs for small users---all in Chapter~12, on compute, algorithm and data supply, with basic research appearing separately in Chapter~8 and without AI attached~\cite{fyp15-ai4s} [P]. Three rules follow for any reader of this literature, and this book applies all three. \emph{Cite the document, not ``the plan.''} \emph{Track phrases across the Recommendations--Outline boundary}, because survival is the test of whether an ambition acquired an owner. And \emph{treat the Outline's silences as data}: a planning system that programs the substrate and omits the epistemology has told you which of the two it intends to be held accountable for. The finding is uncomfortable for this chapter's thesis and is stated in the chapter's own terms rather than buried: on the state's own programming record, Chinese AI for science is an infrastructure commitment with a research aspiration attached, not the reverse.
\end{casestudy}

The money underneath is bench money, distributed widely, and should not be mistaken for a moonshot. Roughly half the project categories in the National Natural Science Foundation's current call foreground AI for science, at individual project budgets of EUR~240{,}000--600{,}000, with topics including generative AI for fusion reactor design and biological component design for extreme-environment biomanufacturing~\cite{nsfc-ai4s} [A]. Two observations follow. The first is that a funding agency restructuring half its call categories around a method is a stronger signal of institutional commitment than any single large programme, because it reaches the whole population of investigators rather than a selected set. The second is that budgets in that range buy graduate students and cluster hours, not frontier training runs---which is consistent with a strategy in which the state supplies the compute centrally (Section~\ref{sec:china-ai4s-infra}) and the grants supply the science.

\section{The output, measured by someone other than China}
\label{sec:china-ai4s-output}

The most useful measurement of Chinese AI-assisted scientific output was produced by the European Commission, which is the reason to use it. Bianchini, Di Girolamo, Ravet and Arranz, publishing on 7~April 2025, examined roughly three million papers across eighty fields over 2000--2022 and found that \textbf{China overtook both the United States and the European Union in AI-assisted scientific output from 2016}, producing more than 25{,}000 AI-assisted papers in 2022 against roughly 15{,}000 for the EU and 12{,}000 for the US, with reported advantages on novelty and citation impact as well as on volume~\cite{ec-ai4s-biblio} [V on the counts, A on the novelty and impact findings]. A second, independently constructed Western measurement points the same way: China leads the Nature Index with output up 22.4\% year on year---the only top-ten country with double-digit growth---and holds the top institution globally along with eight of the top ten in the preceding year's table~\cite{nature-index-2026,nature-index-2025} [V].

Four limits belong immediately beside those numbers, and the first is the one this book has spent a chapter establishing. \emph{Output is not discovery.} Chapter~\ref{ch:prize}'s funnel exists precisely because papers are the first term of a chain whose later terms---validation survival, deployment conversion, lag---are all fractions below one, and a volume lead measures the input side of that chain. \emph{The window ends in 2022}, which is before the model generation this book is about; the study therefore measures the run-up rather than the present state, and its most defensible use is as evidence that the Chinese position was built before the current phase rather than during it. \emph{``AI-assisted'' is a bibliometric classification}, not a causal claim about what the AI contributed. And \emph{a volume lead in a literature with documented quality-control strain} is not automatically a lead in anything: the same chapter records a journal measuring submissions up 42\% and over 30\% of its peer reviews showing detectable AI use. What survives all four qualifications is still substantial, and it is the following: on the two most credible Western-run measurement systems, in the domain this book argues is the largest prize, the leading producer is China, and has been since before the transformer era's commercial phase began.

It is worth stating what would settle the question the bibliometrics cannot, because the absence is itself a finding. A volume lead in AI-assisted publication would be converted into a claim about \emph{discovery} by evidence of the kind Chapter~\ref{ch:prize} required of every Western result it credited: an operational adoption by an agency that did not build the model, a controlled clinical endpoint, an independently replicated synthesis, a proof checked by the community that did not produce it. On that standard the Chinese record is thin in most domains and genuinely strong in exactly one---weather, where third-party evaluation exists in the peer-reviewed literature (Section~\ref{sec:china-ai4s-models})---which is the same domain where the Western record is strongest, and for the same reason: forecasting scores itself daily against observation, and no developer controls the weather. The general instrument for making this comparison does not exist. Nobody operates the production metric of Section~\ref{sec:sci-metric}, in China or anywhere else, which means that the honest current state of knowledge is that we can measure and compare scientific \emph{output} across countries and cannot yet measure or compare AI-assisted scientific \emph{return} in any country. Every claim in this chapter about Chinese leadership is therefore a claim about the input side of Chapter~\ref{ch:prize}'s funnel, and is bounded accordingly.

\section{The institutional layer}
\label{sec:china-ai4s-institutions}

What distinguishes the Chinese institutional layer is not that it exists---every advanced economy now has AI-for-science institutes---but that it was built early, that it is deliberately plural, with several bodies holding overlapping mandates and no single national AI4S laboratory, and that most of these institutions ship artifacts rather than reports. The plurality is step~4 of Chapter~\ref{ch:playbook} in an unexpected setting: the state does not designate the winning laboratory, it funds several and lets the release record select. Table~\ref{tab:china-ai4s} states what each has actually shipped, which is a harder test than what each was founded to do.

\begin{table}[htbp]
\centering\footnotesize
\caption{The Chinese AI-for-science institutional layer as of \datafreeze, scored on shipped artifacts rather than on mandate. Sources:~\cite{ai4s-institutions,cas-scienceone,intern-s1-2025,deepmodeling,pengcheng-cloudbrain,fuxi-weather}.}
\label{tab:china-ai4s}
\setlength{\tabcolsep}{4pt}
\begin{tabular}{@{}p{31mm}p{26mm}p{75mm}@{}}
\toprule
\textbf{Institution} & \textbf{Founded / base} & \textbf{Shipped} \\
\midrule
AI for Science Institute, Beijing & 2021, under CAS academician E~Weinan & Runs the DeepModeling community; DPA potential family; ABACUS density-functional code; DeepFlame combustion platform; AIS-Square model-sharing platform \\
\addlinespace
Shanghai AI Laboratory & Shanghai & The most productive shipper: Intern-S1 (Apache~2.0), Intern-S1-Pro, the Intern Duanyan discovery platform; ten announced 2026 flagship scientific results \\
\addlinespace
Shanghai Academy of AI for Science & 2023, with Fudan & The FuXi weather and climate family, including the Fengshun sub-seasonal model with the National Climate Center \\
\addlinespace
BAAI & Beijing & 2026 ``Wujie'' series: Brain$\mu$ (neuroscience), OpenComplex~2.5 (biomolecules), Physis (physical-world modeling) \\
\addlinespace
Pengcheng Laboratory & Shenzhen & Pengcheng Cloudbrain, the first fully domestic large-scale intelligent computing system; National Science and Technology Progress Award, First Prize \\
\addlinespace
Zhejiang Lab & Hangzhou & Three-Body Computing Constellation: first twelve compute-carrying satellites launched May 2025, stated target $\sim$1{,}000 by around 2030 \\
\addlinespace
CAS ScienceOne & 2025, twelve CAS institutes & Scientific agent platform built on \emph{Chinese open-source foundation models}, integrating AlphaFold and MatterGen as tools; digital-cell platform, collider simulation, telescope coordination \\
\bottomrule
\end{tabular}
\end{table}

ScienceOne deserves two observations that the table cannot carry, and they point in opposite directions. The architectural one: it was built by twelve CAS institutes on Chinese open-source foundation models, and it consumes AlphaFold and MatterGen---the two most celebrated Western scientific models---\emph{as tools underneath its own orchestration layer}~\cite{cas-scienceone} [P]. That is the mechanism of Section~\ref{sec:mirror-move} operating in a scientific setting: whoever holds the layer above treats the layer below as a component, and a Western laboratory that publishes a model without holding the workflow around it has supplied a component to somebody else's system. Its predecessors and successors sit in the same pattern---a national scientific agent platform with access to 170 million papers and 300-plus computing tools, upgraded a year later onto a heterogeneous supercomputing backend, and an autonomous research system wired into a national supercomputing network linking thirty-plus centres~\cite{scienceone,scienceone2,scnet} [P/A]. The epistemic observation cuts the other way and is stated with equal force: ScienceOne's parameters, training compute and open-source status were not disclosed, and its performance claims are self-reported~\cite{cas-scienceone}. A system that cannot be independently evaluated cannot be independently credited, and this book credits its existence and its architecture, not its capability.

What the layer does not contain is as informative as what it does. There is no Chinese equivalent of a national laboratory system with statutory continuity and multi-decade budget lines; the institutions in Table~\ref{tab:china-ai4s} are a mixture of academy institutes, municipally sponsored laboratories and a company, several founded within the last five years, and the plurality that makes them competitive also makes none of them obviously durable. Continuity is the requirement Chapter~\ref{ch:consortium} identifies as decisive for regulated scientific deployment---a clinical model needing two or three years to certify and support five years after that---and it is the one property this layer has not yet had time to demonstrate. That is a real gap in the Chinese position, and it is the mirror image of the West's: the Western system has the institutional continuity and lacks the shipping cadence.

\section{The models, and the fact that matters most about them}
\label{sec:china-ai4s-models}

Four model families carry the Chinese scientific position, and the single fact that matters most about them is not any benchmark number but their licences. This section therefore inverts the usual order of presentation: capability claims are recorded at the grade the evidence supports, which for most of them is developer-reported, and the licence and the training-data composition---the two things that can be checked by anyone with a browser---are treated as the load-bearing facts. That inversion is not a rhetorical device. It follows from Chapter~\ref{ch:prize}'s finding that benchmarks measure components while systems fail, and from Chapter~\ref{ch:trust}'s requirement that a scientific artifact be auditable before it is credited.

\textbf{Intern-S1} is the reference artifact. Released by Shanghai AI Laboratory on 26~July 2025, it carries 241B total parameters---a 235B mixture-of-experts language backbone with a 6B vision encoder---pre-trained on 5T tokens of which \emph{2.5T are scientific-domain}, and it is published under \textbf{Apache~2.0} on Hugging Face, ModelScope and GitHub~\cite{intern-s1-2025} [P]. Its reported scores---ChemBench 83.4, MatBench 75.0, MSEarthMCQ 65.7, MicroVQA 63.9---are self-reported, its comparative claims against Western frontier models are vendor-reported, and no independent replication was located; they are recorded here at that grade and no higher. The same discipline applies to its trillion-parameter successor Intern-S1-Pro of 2~April 2026, whose headline SciReasoner margin over leading closed models is likewise the developer's own measurement~\cite{intern-s1}. The load-bearing facts are not the scores. They are the token mix---half a five-trillion-token pre-training corpus deliberately spent on science, which is a resource-allocation decision no Western laboratory has publicly matched---and the licence.

\textbf{Protenix} makes the licence point unmistakable. ByteDance's AI4Science group released v1.0.0 on 5~February 2026 at 368M parameters and v2 on 8~April 2026 at 464M, both under \textbf{Apache~2.0}, with code, weights, training and inference pipelines and a web server all published; the accompanying claim is that it is the first fully open-source model to outperform AlphaFold3 across diverse benchmark sets at matched training-data cutoff, model scale and inference budget, evaluated over more than six thousand complexes against AlphaFold3, Boltz-1 and Chai-1~\cite{protenix} [P on the release, developer-evaluated on the comparison]. The evaluation was performed with PXMeter, a toolkit the developers themselves built and released---which is better practice than most and still not independent replication, because a harness written by the author of the model under test is a contribution to reproducibility rather than a substitute for it. \emph{The significance is the licence, not the benchmark.} The open-weight ratchet of Section~\ref{sec:maxtier} in Chapter~\ref{ch:model} has reached structural biology, and it reached it from a Chinese laboratory. Note what that does to the position Chapter~\ref{ch:prize} describes: the most diffused AI4S asset on Earth was given away by a Western laboratory as a \emph{database}---two hundred million structures, seventeen million GPU-hours of computation donated---and it now faces a competitor given away as \emph{weights}, which is a different and more consequential act of generosity, because a database is a result and a weight file is a means of production.

\textbf{DeepModeling and DP Technology} are the deepest instance, and the one this book finds hardest to argue with. The stack comprises DeePMD-kit, the deep-potential molecular-dynamics package; \textbf{DPA-2}, published in \emph{npj Computational Materials} in December 2024, covering 73 elements across 18 pre-training datasets spanning alloys, semiconductors, battery materials, drug molecules, ferroelectrics, solid-state electrolytes and catalysis; the ABACUS electronic-structure package; and the OpenLAM initiative launched on 1~January 2024~\cite{deepmodeling} [P]. One detail inside DPA-2 is worth more than its headline: the density-functional labels come from VASP, Gaussian \emph{and} the domestic ABACUS code---which is supply-chain capture executed at the level of the training data, a domestic electronic-structure code inserted into the label pipeline of the flagship domestic atomic model. The community reports more than ten thousand members across more than twenty countries; the commercial arm reports more than three million users across over a thousand universities and research organizations plus corporate clients including PetroChina, CATL and BYD~\cite{deepmodeling,dp-tech} [A]. And the financing settles what kind of entity this is: DP Technology's December 2025 Series~C above RMB~800M (about US\$114M) had \textbf{three of six named investors that are Beijing municipal state funds}~\cite{deepmodeling} [P/A]. This is a state-backed national champion, not a venture outcome, and describing it as a start-up misstates the instrument. Code, weights and datasets are open. The sentence this book is obliged to print, because it is what the evidence located actually supports and because it cuts against the framing a Western reader will bring: \emph{on the evidence available at \datafreeze, this is among the most genuinely open scientific-AI ecosystems anywhere, Western or Chinese.}

\textbf{Weather is the one domain with genuine external validation, and it inverts the openness picture.} Huawei's Pangu-Weather has been checked by third parties in the peer-reviewed literature---evaluations in \emph{Remote Sensing}, \emph{Meteorological Applications} and \emph{Weather} during 2025, the last covering FengWu as well---which makes this the only Chinese AI4S domain where the capability claim rests on somebody other than the developer~\cite{fuxi-weather} [V]. The FuXi family, by contrast, is strong on deployment and weak on verification: China's first purely data-driven end-to-end forecasting system, ingesting satellite remote sensing directly, operationally deployed across Chinese meteorological departments and integrated into an early-warning platform serving forty-three countries, selected among China's 2025 Top Ten Meteorological Science and Technology Advances in March 2026---while its headline claims, medium-range ensembles beating ECMWF on 98.1\% of variables by CRPS and sixty-day forecasts in three minutes, are developer-reported and not independently replicated, and the open-weight status of the current system is not stated~\cite{fuxi-weather} [P on deployment, developer-reported on skill]. The meteorological administration's own operational family, one member of which has been open-sourced, sits alongside~\cite{cma-ai}. Structural biology and weather thus form an instructive pair: the domain where the Chinese artifact is most open is the one where the claim is least externally checked, and the domain where the claim is best checked is the one where the artifact's licence is unstated.

\textbf{And now the counter-observation, which must be made and made prominently.} The newest and most capable Chinese scientific platform---Shanghai AI Laboratory's Intern Duanyan, announced 17~July 2026, a five-layer discovery stack with more than two hundred cross-disciplinary agents across six domains and a dry-lab/wet-lab closed loop---has \textbf{no stated open-source release}, and is access-by-invitation~\cite{intern-s1-2025} [P]. Read against the mechanism of Chapter~\ref{ch:model}, this is not an anomaly but a prediction coming true. The ratchet compels \emph{weights} onto the commons because adoption is the currency and the first defector loses the derivative ecosystem; it compels nothing whatever about a \emph{platform}, because a closed-loop agentic laboratory with wet-lab integration is not a file anyone can fork, and its rent lives in the orchestration and the instrument access rather than in the parameters. The openness documented in this section may therefore be a phase rather than a principle---exactly what Chapter~\ref{ch:playbook}'s step~6 predicts of any layer once it starts producing rent, and exactly the policy-and-continuity risk that Chapter~\ref{ch:consortium} names first among the three its insurance is meant to cover. That risk is no longer hypothetical. It is instantiated, and it is instantiated at the top of the Chinese capability stack.

\section{The infrastructure, and the prize the state actually gave}
\label{sec:china-ai4s-infra}

Pengcheng Cloudbrain is the first fully domestic large-scale intelligent computing system, built by Pengcheng Laboratory with Huawei, Peking University and Tsinghua, and on 8~July 2026 it was announced as the winner of the \textbf{National Science and Technology Progress Award, First Prize} for the 2025 cycle; it is reported to support training for more than half of China's domestic foundation models and to have guided the construction of twenty-five intelligent computing centres~\cite{pengcheng-cloudbrain} [P on the award, A on both ratios]. \emph{The state's top science prize went to the machine, not to a discovery made on it.} A national prize list is a revealed preference over what a state believes is scarce, and this one says that the scarce thing is not a result but the domestic capacity to produce results---which is the argument of Chapter~\ref{ch:compute} restated by the awarding body rather than by this author, and the reason Chapter~\ref{ch:sizing} measures capacity rather than publications.

The 15th Five-Year Plan Outline supplies the instrument underneath. Its compute provisions direct the study and construction of \emph{ultra-large-scale intelligent computing clusters}; they name government \emph{purchase} and \emph{leasing} of compute as explicit policy instruments; and they set the aim of making compute ``universally accessible and easy to use'' while lowering costs for smaller users~\cite{fyp15-ai4s} [P]. That is move~M2 of Chapter~\ref{ch:playbook}---demand before supply---applied to research computing rather than to solar modules: the state does not merely build the machine, it commits to being its customer, and it prices access down for precisely the long tail that Chapter~\ref{ch:consortium} identifies as the population any scientific-compute institution must serve. Note also which of the two documents of Section~\ref{sec:china-ai4s-policy} this provision lives in. The paradigm-transformation sentence dropped out between the Recommendations and the Outline; the compute-purchase provision survived. The state plan programs the substrate, not the epistemology, and the prize list agrees with the plan.

The inversion is worth putting to Western science ministries directly, because it is the most transferable thing in this chapter and it costs nothing to act on. A national prize system is one of the cheapest instruments a government holds and one of the most legible signals it sends: it tells a generation of researchers which career produces recognition. Western prize lists in this domain decorate results---the structure, the forecast, the molecule---and treat the infrastructure that produced them as procurement. The Chinese list, in its 2025 cycle, decorated the infrastructure. Neither choice is obviously right; a state that only rewards machines will get machines and no science, and a state that only rewards discoveries will find, a decade later, that it rents the means of producing them. But the asymmetry is real, it is recent, and it points the same way as every other finding in this chapter: toward a programme organized around productive capacity rather than around outputs, which is precisely the organizing principle Part~I found in solar, in batteries and in vehicles.

One quantity this chapter declines to supply. There is no public figure for how much compute the Chinese research system actually receives, in any unit comparable to the BGPU accounting of Chapter~\ref{ch:sizing}, and no Chinese instantiation of that chapter's national demand model exists. This chapter therefore does not place China on the sizing ladder, does not estimate its per-researcher provision, and does not assert that its research population is better served than Japan's or America's---all three of which would be easy to write and none of which the evidence supports. What can be stated is the instrument and its direction: a state that buys and leases compute on behalf of its researchers and drives the price down for small users is acting directly on the gap that Chapter~\ref{ch:sizing} measures at roughly tenfold in the countries where it has been measured, whatever level it in fact reaches.

\section{The programme, scored on the playbook}
\label{sec:china-ai4s-playbook}

Chapter~\ref{ch:playbook} distilled six strategic steps and nine tactical moves, and the discipline this book imposes on itself is to score campaigns against them rather than to admire them. The AI4S programme scores as follows.

\textbf{Step~1, target selection (M1).} The discontinuity chosen is not a product category but the research paradigm itself---``use artificial intelligence to lead a transformation of the scientific research paradigm'' is a target-selection sentence in the exact form of Chapter~\ref{ch:playbook}'s first step~\cite{fyp15-ai4s}. What is stranded if the paradigm changes is the incumbent's accumulated advantage: the seat-licensed scientific software stack, the closed frontier platforms, the instrument and publication ecosystem, and the tacit practice of research groups organized around them. State designation is present, dated and layered---the 2023 ministerial deployment, the 2025 State Council Opinions, the 2025 Party Recommendations---though, as Section~\ref{sec:china-ai4s-policy} established, thinning as it moves from Party ambition to state programming.

\textbf{Step~2, demand before supply (M2).} Government purchase and leasing of compute, subsidized access, the universal-accessibility language, and twenty-five guided intelligent computing centres~\cite{fyp15-ai4s,pengcheng-cloudbrain}. This is structurally the feed-in tariff of the solar column and the Top Runner procurement mechanism, redirected at research capacity: guaranteed demand for domestic machines, delivered as subsidized supply to a scientific user base that could not otherwise afford it.

\textbf{Step~3, overcapacity as strategy (M3): partially, and the evidence is weaker than the pattern.} Twenty-five guided intelligent computing centres and an Outline provision directing the study and construction of ultra-large-scale clusters have the shape of the solar and battery build-outs---capacity funded ahead of demonstrated demand, with the losses absorbed locally and the cost curve retained nationally~\cite{pengcheng-cloudbrain,fyp15-ai4s}. But the shape is not the finding. No public utilization figures for those centres were located, and without them the difference between strategic overbuild and stranded provincial assets cannot be established from open sources; this book records the instrument and declines to score the step [S]. It is worth naming what evidence would settle it: published occupancy or allocation data from the intelligent computing centres, of the kind Chapter~\ref{ch:sizing}'s accounting unit was designed to make comparable across countries.

\textbf{Step~4, selection rather than planning (M4).} Six-plus institutions with overlapping mandates, two competing weather families, several scientific-model families, and no single designated national AI4S laboratory (Table~\ref{tab:china-ai4s}). Chapter~\ref{ch:consortium} draws precisely this lesson from the Chinese case---that the government's role was to set the stage and permit intense bottom-up activity---and it is worth noting that the lesson is drawn there from a working discussion among Western centre directors, not from Chinese sources, which makes it a convergent reading rather than a borrowed one.

\textbf{Step~5, commoditize the adversary's rent layer (M6).} There are \emph{two} rent layers here, and the Western debate has priced only one. The first is the closed frontier platform, the subject of Chapter~\ref{ch:model} and of the arithmetic in Section~\ref{sec:ai4s-econ}, where a 20--60$\times$ purchased-token premium converts a one-seventh workload share into a three-quarters-to-nine-tenths budget share. The second is the one nobody in the AI debate mentions: the \emph{subscription-priced scientific software stack}---commercial density-functional and molecular-dynamics codes, cheminformatics suites, simulation packages licensed per seat per year, structure and sequence resources behind institutional paywalls---which has collected rent from the world's researchers for three decades. Open scientific weights, open codes and open datasets price both layers at marginal cost simultaneously. Section~\ref{sec:mirror-move}'s result governs the consequence: the rent does not vanish, it relocates, and it relocates toward compute, toward platforms with instrument access, and toward whoever holds the orchestration layer---which is what Intern Duanyan and ScienceOne are.

\textbf{M5, supply-chain capture.} The chain is visible end to end: domestic open foundation models under ScienceOne rather than Western ones; ABACUS as a domestic density-functional code inside the label pipeline of the flagship domestic atomic model; domestic silicon underneath, supplied through the Huawei partnership that built the first fully domestic large-scale intelligent computing system---the machine the state chose to decorate [P on the partnership and the award, A on the accelerator line]; and, at the newest link, compute-carrying satellites at Zhejiang Lab against a stated target of roughly a thousand by around 2030~\cite{ai4s-institutions,deepmodeling,pengcheng-cloudbrain}. Model, framework, code, chip, machine---the same sequence as Chapter~\ref{ch:playbook}'s AI row, run through a scientific application rather than a commercial one.

\textbf{M8, standards capture: not achieved, and available.} The one visible instance is the potential-energy-model interface, where DPA and OpenLAM are becoming a de facto interface for machine-learning interatomic potentials in the way file formats become standards---through a ten-thousand-member community across twenty-plus countries rather than through a committee~\cite{deepmodeling}. Beyond that there is no Chinese scientific-model certification scheme, no AI4S standards body, and no scientific equivalent of MLPerf. That vacancy is exactly the one Section~\ref{sec:sci-metric} of Chapter~\ref{ch:prize} identified and proposed a metric for, and it remains unfilled by anybody---which is a strategic opportunity for the institution of Chapter~\ref{ch:consortium} and a reason Chapter~\ref{ch:trust}'s certification argument is urgent rather than academic.

\textbf{The industrial hinge, which is this chapter's distinctive finding.} Of the ten joint results Shanghai AI Laboratory announced at the 2026 World AI Conference, three are aerospace design with a state aerospace research institute and a commercial rocket firm, electroplating-materials formulation with Xiamen University, and an analogue chip-design agent with Shanghai Jiao Tong University claiming five- to tenfold sizing-optimization efficiency and a halving of layout-planning cycles~\cite{ai4s-industrial} [P on the announcement, developer-reported on the efficiency figures]. A ``scientific'' programme delivering into aerospace, industrial chemistry and semiconductor design is not a scientific programme in the Western sense of the phrase; and the analogue-design case is the flywheel of Section~\ref{sec:flywheel} pointed directly at the layer Chapter~\ref{ch:semi} identifies as the binding physical constraint. The policy text makes the intent textual rather than inferred: the science pillar's second sub-clause chains AI-driven research to ``engineering realisation and product landing''~\cite{ai-plus-science}. That is the sense in which this chapter's title is a description rather than a characterization.

\textbf{And the detail that cuts against the framing, printed because it does.} Three of those same ten results carry foreign co-authors---Columbia, Cambridge, University College London~\cite{ai4s-industrial} [P]. The programme is not hermetically sealed, and any account that treats Chinese AI for science as a closed national system is wrong on the byline evidence. This matters directly for the next section, because an institution designed to exclude Chinese participation would be excluding collaborators that Western universities already have, on projects already published.

\textbf{Step~6 and M9, cartelize and gatekeep: not observed.} There is at \datafreeze no restriction on Chinese scientific model release, no export control on scientific weights, and no gatekeeping behaviour in this domain of the kind rare earths and polysilicon supply. The single datapoint consistent with the endgame is the unstated licence on the most capable platform, and one datapoint is a flag rather than a finding [S].

\section{What this does to the consortium proposal, and what it does not}
\label{sec:china-ai4s-consortium}

If this chapter is right, then the case usually offered for a Western or neutral open scientific-AI consortium is not merely weak but false, and it should be abandoned rather than defended. That case says: the scientific commons is thin, closed platforms dominate, and public money must create an open base. On the evidence of Section~\ref{sec:china-ai4s-models}, the open base exists. It comprises Apache~2.0 scientific foundation models at 241B and trillion-parameter scale, an openly licensed structure-prediction model with published code, weights and pipelines, and an atomic-potential ecosystem whose code, weights and datasets are all open and whose community spans twenty-plus countries. A consortium justified as \emph{creating} the commons would be arguing against a fact, and would be found out quickly by the researchers it proposed to serve, who are already downloading the artifacts.

The defensible case is different, narrower, and stronger, and it is the one Chapter~\ref{ch:consortium} actually makes: \emph{the commons has one principal contributor and no governance the rest of the world is party to.} Each of the three risks that chapter names is sharpened rather than softened by the evidence assembled here. \emph{Continuity}: the most capable Chinese scientific platform has no stated open release, which converts the argument from a hypothetical about future Chinese policy into an observation about present Chinese practice. \emph{Verification}: ScienceOne's parameters and training compute are undisclosed, FuXi's headline skill claims are developer-reported, Intern-S1's comparative benchmarks are vendor-reported, and Protenix---the most rigorous of the four---is evaluated by a harness its own authors wrote. Four of the most consequential artifacts in the field cannot be checked in the way Chapter~\ref{ch:trust} requires, and the correct response to that is an evaluation institution, not a rival model. \emph{Monoculture}: if a single national ecosystem supplies the interatomic potentials, the scientific foundation models and the agent frameworks that a generation of published results depends on, then the systemic argument of Chapter~\ref{ch:order66} applies to the scientific literature itself, and applies whether or not anyone has behaved badly.

The strongest version of the opposing case must also be printed at full strength, because it is genuinely strong. A genuinely open scientific commons is a public good regardless of who funds it: a researcher in Nairobi or Jakarta running an open atomic potential on open weights is better off than one who cannot, and the nationality of the contributor does not enter that welfare calculation anywhere. A consortium built to \emph{exclude} its largest contributor would be a worse instrument than one built to include it---it would forfeit the artifacts, it would forfeit the neutrality that Table~\ref{tab:consortium-spec} identifies as its only durable asset, and it would substitute an origin rule for the evidence-based evaluation that is the sole product an evaluation body has to sell. And there is a harder form of the objection still. If AI for science is the largest prize, and scientific knowledge is non-rival, then a policy that slows the commons in order to change its ownership destroys more value than it redistributes, and does so in the one domain where Chapter~\ref{ch:prize} calculated the social return at above four dollars on the dollar. On that reading the correct Western response to Section~\ref{sec:china-ai4s-models} is not a rival consortium at all but a larger contribution to the same commons, plus the evaluation apparatus nobody yet operates.

One asymmetry is worth isolating, because it is the pivot on which the disagreement actually turns and it is not about models at all. Contribution to a commons and governance of a commons are different goods, and the second does not follow from the first. The Chinese programme supplies artifacts at a rate no other single actor matches; it does not supply---and has not been asked by anyone to supply---an evaluation regime, a revocation procedure, a provenance standard, a certification ladder, or a forum in which a non-Chinese institution has standing to contest a result. Those are the five functions Chapter~\ref{ch:consortium} specifies operationally and the ones Chapter~\ref{ch:trust} argues nobody currently owns. Framed that way the choice put to the rest of the world is not between a Chinese commons and a Western one; it is between a commons with artifacts and no governance and a commons with both. That framing has the additional merit of being buildable: an evaluation and certification institution can be constituted, funded and staffed at a small fraction of what a rival model programme costs, and unlike a rival model programme it becomes more valuable, not less, as the artifacts it evaluates get better---whoever makes them.

This chapter does not resolve that tension; Chapter~\ref{ch:consortium} does, and does so after Chapter~\ref{ch:sizing} has supplied the missing quantity. What this chapter has established is the shape of the question those two chapters inherit. The prize is the largest one on the table. The most industrially organized programme aimed at it is Chinese, is legible in primary policy text, chains discovery to product by the wording of that text, runs on a substrate the state's own prize list ranks above the discoveries made on it---and distributes a substantial part of its output for free, under permissive licences, to everyone including its competitors. The question left for the rest of Part~IV is therefore not whether to build on that commons, which every scientist reading this book is already doing, but whether an open commons with one principal contributor and no shared governance is a commons or a dependency, and what instrument short of exclusion changes the answer.
